% §3.2 OR and NOT
%
% Cross-reference resolution status:
%
% RESOLVED:
%   \Cref{sec:and-under-noise}, \Cref{eq:and-fnr},
%   \Cref{eq:and-fpr-general}     -- chapter 3 §3.1
%   \Cref{sec:composition-theorem}-- chapter 3 §3.3
%   \Cref{def:bernoulli-bool}, \Cref{ax:element-indep} -- chapter 4
%
% DEFINED here: \label{sec:or-not}, \label{eq:not-Q},
%   \label{eq:or-fpr-direct}, \label{prop:boolean-closure}
%
% Monograph register: NOT swaps the rate pair; OR via De Morgan and via
% the truth table agree; the AND/OR duality; closure of Bernoulli[Bool]
% under Boolean algebra.

\section{OR and NOT}
\label{sec:or-not}

NOT is immediate and OR follows from it by De Morgan, completing the Boolean
basis.

\subsection*{NOT}

Complementing the output of a $\Bernoulli{\BoolType}$ instance turns each false
negative into a false positive and conversely: a $\bot$ reported on a true
$\top$ becomes a $\top$ reported on the complemented true $\bot$. So NOT swaps
the rate pair, $(\fprate, \fnrate) \mapsto (\fnrate, \fprate)$, which on the
channel matrix is the row exchange
\begin{equation}
\label{eq:not-Q}
  Q_{\neg} \;=\;
  \begin{pmatrix}
    1 - \fprate & \fprate \\
    \fnrate     & 1 - \fnrate
  \end{pmatrix}.
\end{equation}
NOT is an involution on $\Bernoulli{\BoolType}$: applying it twice restores the
rates, the noisy-bit form of $\neg\neg b = b$ (at the level of the channel
matrix, not of any realized output).

\subsection*{OR}

By De Morgan, $\tilde{b}_1 \vee \tilde{b}_2 = \neg(\neg \tilde{b}_1 \wedge \neg
\tilde{b}_2)$. Negating the inputs swaps their rates, the AND of the negated
instances has false-negative rate $\fprate_1 + \fprate_2 - \fprate_1 \fprate_2$
by \eqref{eq:and-fnr} with rates swapped, and the outer NOT swaps that into the
OR's false-positive rate. The truth table confirms it directly: the observed
OR is $\bot$ only when both instances report $\bot$, so a false positive (true
$\bot\bot$, observed $\top$) is
\begin{equation}
\label{eq:or-fpr-direct}
  \fprate_\vee \;=\; 1 - (1 - \fprate_1)(1 - \fprate_2)
               \;=\; \fprate_1 + \fprate_2 - \fprate_1 \fprate_2,
\end{equation}
the inclusion-exclusion probability that at least one instance false-positives.
This is prior-free. The OR's false-negative rate, dually, is prior-dependent
and equal to the product $\fnrate_1 \fnrate_2$ in the all-positive regime.

The duality is exact: AND has a prior-free false-negative rate
\eqref{eq:and-fnr} and a prior-dependent false-positive rate, while OR has a
prior-free false-positive rate \eqref{eq:or-fpr-direct} and a prior-dependent
false-negative rate. AND suppresses false positives and inflates false
negatives; OR does the reverse; NOT exchanges the two regimes. In the
symmetric case $\fprate_i = \fnrate_i = \fprate$, both inflate by the same
$2\fprate - \fprate^2$, and De Morgan's duality is manifest.

\begin{proposition}[Closure under Boolean operations]
\label{prop:boolean-closure}
If $\tilde{b}_1, \ldots, \tilde{b}_n$ are $\Bernoulli{\BoolType}$ instances,
every Boolean formula built from them with AND, OR, and NOT is a
$\Bernoulli{\BoolType}$ instance whose rates are polynomials in the input
rates.
\end{proposition}

\begin{proof}
Induct on the formula. NOT preserves $\Bernoulli{\BoolType}$ by
\eqref{eq:not-Q}; AND does by \eqref{eq:and-fnr} and
\eqref{eq:and-fpr-general}; OR does by De Morgan from NOT and AND. Since
$\{\wedge, \vee, \neg\}$ is a complete basis and every formula is a finite
composition of these, the induction closes.
\end{proof}

Closure is what makes the framework composable: a Boolean circuit of noisy
gates has, at every stage, a $\Bernoulli{\BoolType}$ output whose error
propagates in closed form, with the prior-dependence of the AND and OR rates
inherited as in \Cref{sec:and-under-noise}. The next section isolates the case
where that propagation is cleanest, the chained one-sided gate, and states it
as the composition theorem.
